\section{Existence of sequential equilibrium}\label{sec:existence}

\begin{theorem}\label{thm:existence}
In every Krusell--Smith economy, from every aggregate shock $z_0$ and every initial distribution
$\mu_0$, a sequential equilibrium exists, with aggregate labour equal to the mean efficiency at
every node.
\end{theorem}
\formal{exists\_sequentialEquilibrium, exists\_sequentialEquilibrium\_labour}

\noindent
The unknown is a path of capital on the tree of shock histories, with values in $[0,\abar]$. These
paths form a compact convex subset of a locally convex space in the product topology. The household
value function is continuous in the path: a depth-weighted sup norm makes the Bellman operator a
contraction with a modulus uniform in the path, and the weights make the tree's distant nodes
negligible. Strict concavity gives a unique policy, continuous in the path by Berge's theorem; the
distributions on the tree are then weakly continuous in the path by induction on depth, and so is
mean capital. Schauder--Tychonoff gives a fixed point.

\begin{theorem}\label{thm:existenceInada}
The same holds with utility defined only on $(0,\infty)$ and marginal utility tending to infinity at
zero --- logarithmic and CRRA utility --- provided labour income is bounded below by a positive
floor; optimal consumption never reaches the consumption floor used in the construction.
\end{theorem}
\formal{KSInada.exists\_sequentialEquilibrium}

\begin{theorem}\label{thm:existenceCD}
With logarithmic utility and Cobb--Douglas prices, a sequential equilibrium exists, and prices are
the true marginal products wherever capital lies in the clamping box. With CRRA utility, an explicit
lower end $\underline K$ of the capital box (depending on $\beta$, $\gamma$, the persistence of the
aggregate chain and the technology) keeps every equilibrium's capital above $\underline K$: the
clamp never binds.
\end{theorem}
\formal{logEconomy\_exists\_equilibrium, crra\_exists\_unclamped\_equilibrium}

\noindent
The lower bound adapts \citet[Lemma 9]{cao2020}: a per-household Euler inequality, aggregated with
Jensen's inequality for the convex function $c\mapsto c^{-\gamma}$, bounds consumption today by
expected consumption at the child nodes, and an induction on the tree pushes capital above
$\underline K$.
